\egasection{§10\quad 超平面截面束与簇的爆破的纤维化}\source{法稿 \texttt{EGA5\_francais.pdf} 物理第56--60页上部（手写页码53--57）；英译 \texttt{EGAVaiello.pdf} 第38--41页；第十九版法英对读}
\egap{10.1}
设 $Z$ 是 $P^\vee$ 中相对于一个可构造性质 $\mathcal P$ 的例外集，并假设 $Z$ 为可构造子集，$S=\Spec(k)$。例外集的余维按其闭包的余维理解。以下关于直线的判据以 $\dim P^\vee\geq1$ 为前提。\transnote{原稿直接讨论直线，未单列零维参数空间；零维射影空间不含直线，不能使用这里的直线存在性判据。}我们将看到（参见第12节，那里将补述一些无疑本应在前几节出现的内容），条件
\[
\codim(\overline Z,P^\vee)\geq2
\]
等价于：几何意义下充分一般的直线 $L$ 不与 $\overline Z$ 相交。对于闭集 $\overline Z$，只要存在一条与之不相交的几何直线便足够。\cmt{D19-R04}{法稿第56页在“或者 $\overline Z$”及最后无限域的判据中均可辨上划线；旧译只对英译，因而把两处一概记为“没有可辨的上划线”，现予更正。不过法稿“只需存在一条直线”这句的末字本身误打为 $k$；英译改印 $Z$，并附 P.B. 的“应为 $Z$；A.G.的错误”批语，未清楚保留闭包。正文统一为 $\overline Z$：单条直线避开一般可构造集 $Z$ 并不足够。反例是射影平面中一条直线删去一点；过所删点的另一条直线避开该集合，而其闭包余维仍为1。英译原脚注55：“请询问 A.G.”。}
若 $k$ 无限，上述余维条件又等价于存在一条定义于 $k$、不与 $\overline Z$ 相交的直线 $L\subset P^\vee$。\transnote{法稿此句为 $\overline Z$，英译印为小写 $\bar z$；正文恢复法稿。闭集余维至少2时，避开它的直线构成 Grassmann 簇的非空开集；其仿射坐标图在无限域上有稠密的 $k$-有理点。反之，任何余维1的射影闭子集都与每条几何直线相交。这里区分几何直线与定义于 $k$ 的直线。}
我们把 $L$-预概形 $Y_L$ 称为由 $P^\vee$ 中直线 $L$ 定义的 $X$ 的一个超平面截面线性束（此定义对任意 $S$ 都有效）。

于是，将上述考虑与第8、9节的结果结合，首先在 $S$ 为一个无限基域的谱的情形下，便得到这样的束的存在性判据：所有纤维 $Y_\xi$（$\xi\in L$）都满足性质 $\mathcal P$。考虑到8.2，若对 $X$ 的每一个伴随素循环 $T$ 都有
\[
\dim f(T)>0,
\]
则可以要求束 $Y_L$ 在 $L$ 上平坦，办法是取性质
\[
\mathcal P'=\mathcal P+\text{$\phi_\xi$ 的正则性条件}.
\]
在 $S$ 任意的情形下，若 $k(s)$ 无限，而且已知 $Z$ 为闭集，则仍可按照7.1的方法，在 $S$ 的给定点 $s$ 的某个开邻域上构造这样的束。（第5节的结果以及 $X\to S$ 固有的假设，在各种不同情形下保证上述闭性。）

为了妥善处理此事，宜在作出这种一般性说明后，给出若干汇总性陈述，将前面的结果实际应用于若干这类性质（也包括模的性质）。至少，必须在这里用线性束的语言重新表述9.5；这是一项在几何应用中经常使用的事实。\transnote{此句是原稿的写作计划及评语，并不意味着底本中缺失的9.5已经补齐。}

% DOT_IMPORT 16 10.2_校订续译_工程
\egap{10.2}
借助极性对应，$P^\vee$ 中的一条直线 $L$ 对应于 $P$ 中余维为2的线性子簇 $L^0$（$S$ 任意）。置\transnote{法稿第57页完整写作 $T=X\times_P L^0$。英译第39页把 $X$ 排成小写 $x$ 且缺少纤维积的清楚排法；本轮据法稿恢复，不再将其称为无法辨认的原稿缺文。}
\[
T=X\times_P L^0.
\]

还可以如下描述 $T$：$L$ 由 $E^\vee$ 的一个秩为2的局部自由商模定义；等价地，由 $E$ 的一个处处秩为2、局部为直和因子的子模 $F$ 定义。考虑复合同态
\[
F_X\longrightarrow E_X\longrightarrow\mathcal O_X(1),
\]
则 $T$ 正是这个复合同态的零点概形；换言之，它由理想 $J$ 定义，其中 $J$ 是相应同态（通过张量乘以 $\mathcal O_X(-1)$ 得到）
\[
F_X(-1)\longrightarrow\mathcal O_X
\]
的像。

在本节用于有限型簇的局部诺特情形，假设这个同态是正则的：在零点概形 $T$ 的每一点附近取 $F_X(-1)$ 的一个全序局部基，其像构成 $\mathcal O_X$-正则序列；在 $X\setminus T$ 上不再要求一个生成单位理想的序列满足“真理想”的约定。这个条件与所选局部基无关，并可内在表述为：
\[
F_X(-1)\otimes_{\mathcal O_X}(\mathcal O_X/J)
\longrightarrow J/J^2
\]
为同构，而且 $V(J)=T\to X$ 为正则浸入。\cmt{D19-R05}{法稿第57页本来以任意预概形为背景，原句为：“假设这个同态是正则的，意思是在局部写出 $F_X(-1)$ 的一个全序基，其像构成正则序列。这个条件与基的选择无关，且也可以内在地说成……”。本轮保留该内在同态（法稿可辨，英译漏写），同时把已核实的等价与换基不变性限定在局部诺特情形。一般非诺特环中，普通正则列、Koszul正则列等概念必须区分，不能无说明地套用换序、换基与这些等价。此处限定是本轮审证的范围，不把未审的一般情形宣称为假；原稿的任意预概形设定仍记录在此。空中心按零点概形附近的约定允许，正是10.4所需。}
〔原稿附记：应当在某处阐明一般情形，即预概形 $X$ 上有一个同态 $G\to\mathcal O_X$，其中 $G$ 局部自由；例如可放在关于正则浸入的一节。〕于是有：

\par\noindent\textbf{定理（10.2）。}
在上述已明确的正则性假设下，带有典范投影 $Y_L\to X$ 的线性束 $Y_L$，以唯一的方式在 $X$ 上同构于预概形 $X$ 沿中心 $T$ 的爆破。

为了理解这个定理的意义，宜在本节或本条开头指出：若 $S=\operatorname{Spec}(k)$，则对 $P^\vee$ 中一条“充分一般”的直线 $L$，上述正则性条件成立（参见第12节所指的补述，即关于5.3的补述）。\transnote{这是原稿对后文的预告，不应理解为无须约定定义域便可在任意有限域上选取一条所需的 $k$-有理直线；其精确范围仍须结合5.3与第12节核对。}
在下文构造本条开头所指出或预期的“优良”线性束时，我们可以要求所述束满足这一条件（它与下述条件属于同一类型，但并不相同：对每个 $\xi\in L$，要求 $\phi_\xi$ 是 $X_{k(\xi)}$-正则的）。应把所讨论的条件纳入拟作的汇总性陈述。

另一方面，实际上10.2主要用于9.5的情形，因此，最好在10.2之后才用束的语言重新表述9.5，以便在该陈述中同时纳入束与一个爆破的同构（即给出便于适当引用的情形描述）。按原稿这里拟议的应用，对于无限域 $k$ 上每一个维数 $\ge2$ 的光滑、几何连通的射影概形 $X$，我们便得到一种方法，找出一个非空的光滑闭子概形，它在其每一点的余维均为2，使得爆破后的概形具有到 $\mathbf P^1$ 的一个纤维化：所有纤维均几何整，并且除至多有限个纤维外，其余纤维都光滑；那些例外纤维至多有一个几何奇异点，而且这样的点在相应参数点 $\xi$ 的剩余域 $k(\xi)$ 上有理，并且在几何意义下是普通奇点。\transnote{法稿第58页和英译均写“在 $k$ 上有理”。纤维在参数点 $\xi$ 的剩余域上定义，故相应唯一几何奇异点的有理性应表述为在 $k(\xi)$ 上有理，不能无条件说成在 $k$ 上有理；正文据此订正其基域。这整段仍是原稿拟议的应用，所引用的（9.5）在现有底本中残缺，译校不声称已补齐它或独立证明全部存在性。}

这说明，深入研究这种带奇异纤维的纤维化是重要的（这种研究在目前才刚刚开始）：它可在一定程度上把维数为 $d$ 的光滑射影簇的研究，归约为对维数为 $d-1$、容许有普通奇点的射影簇的研究（这些簇组成依赖一个参数的族）。\transnote{法稿第58页的手写插入可辨为“pouvant avoir des singularités ordinaires”（容许有普通奇点）；正文据此恢复英译中语法不顺的“or”，并整理括号。前一段的“主要”也据法稿“utilisé surtout”恢复；旧中文“只用于”过强。}

陈述10.2或多或少是下面命题的直接推论；后者完全独立于超平面截面这一主题，无疑更适合放在另设的一段“正则浸入”中。\footnote{原脚注56：“请询问A.G.”}

〔底本编排批注〕命题10.3被划去了。请询问A.G.，这是否是他的本意。

% DOT_IMPORT 17 10.3_校订续译_工程
\egas{10.3}{命题（底本标记已划去；以下保存其文字）}
设 $X$ 为预概形，$G$ 为 $X$ 上的拟凝聚模，
\[
u:G\longrightarrow\mathcal O_X
\]
为一个同态，并置 $J=u(G)$、$T=V(J)$。设 $\widetilde X$ 是由 $X$ 沿 $T$ 爆破而得的预概形。\transnote{法稿第59页以带波浪号的 $\widetilde X$ 区分爆破后的预概形；英译漏掉波浪号。本轮据法稿恢复，不把它归为原作者混用对象。}
另一方面，考虑 $P=\mathbf P_X(G)$、典范同态
\[
G_P\longrightarrow\mathcal O_P(1)
\]
及其核 $H$，于是有正合列
\[
0\longrightarrow H\longrightarrow G_P
\longrightarrow\mathcal O_P(1)\longrightarrow0.
\]
再考虑同态 $u_P:G_P\to\mathcal O_P$，以及拟凝聚理想\transnote{法稿第59页为大写 $P=\mathbf P(G)$；英译误作小写 $p=p(G)$，正文恢复大写并明示底空间 $X$。$G_P$ 表示 $G$ 在 $P$ 上的拉回；$u_P$ 为相应拉回同态。底本的“quadi-coherent”按“quasi-coherent”译出。}
\[
K=u_P(H)\subseteq\mathcal O_P.
\]

则 $\widetilde X$ 典范地同构于 $V(K)$ 的一个闭子概形。

若进一步 $G$ 局部自由且秩有限，$X$ 局部诺特，并且 $u$ 是上述意义下正则的，则这个同构就是 $\widetilde X$ 与 $V(K)$ 本身的同构。在这种情形下，$H$ 为 $P$ 上的局部自由模，而同态
\[
H\longrightarrow\mathcal O_P
\]
（其零点预概形为 $\widetilde X$）也是正则的。\cmt{D19-R06}{“零点预概形”据法稿第59页清楚可辨的“préschéma des zéros”恢复；英译此处为“???”，旧译因未见法稿而留空。本段在法稿中与10.3全条一起被大叉划去。本轮完整保留这份被删除的草稿，但不改变其身份。原文只写“$G$ 局部自由且 $u$ 正则”，未另列秩有限或局部诺特；正文明确本轮给出局部证明的范围，一般性问题仍留待统一正则性的约定。前一段关于典范闭浸入的断言仍保留其任意拟凝聚 $G$ 的范围。}

\subsection*{原稿紧随其后的论证说明}
第一个断言几乎是平凡的。第二个断言是一个不会造成困难的练习（我没有详细做过，心想你和我一样能够把它推导出来）。\transnote{“几乎是平凡的”及括号内的第一人称均为原稿文字，并非本译者对证明充分性的判断。现有底本此处没有详细证明，不能将这一说明当作已完成的独立审证。}

% DOT_IMPORT 18 10.4_校订续译_工程
\subsection*{（10.4）之前的说明}
若在10.2中取 $S=\operatorname{Spec}k$，且 $X$ 的维数 $\le1$，则在上述局部诺特及空中心约定下，所作正则性假设等价于 $T=\varnothing$，因而 $Y_L\to X$ 为同构。\cmt{D19-R07}{法稿第59页明确写 $\dim X\le1$；英译误成 $\ge1$。旧译为避免猜测，只写下述光滑曲线情形，本轮据实见法稿恢复“至多1维”。非空中心处若有长度2的正则列，则局部环深度至少2；在维数至多1的诺特局部环中不可能。反向依10.2在零点概形附近定义正则性的约定，空中心时条件为空。这里既不是由数学常识猜回不等号，也不把结论扩大到未约定正则性的任意非诺特预概形。}因此，将此与9.2结合，得到：
\egas{10.4}{推论（（10.2）的推论）}
设 $X$ 是无限域 $k$ 上射影空间 $P$ 中的一条光滑、几何连通的射影曲线（取其给定闭浸入），且 $n\geq2$。则存在 $P$ 上的一个 $n$ 次形式的线性束，它定义一个态射
\[
X\longrightarrow\mathbf P^1,
\]
并具有下述性质：这个态射泛平展，次数为 $d$；对于 $\mathbf P^1$ 的每一个几何点 $s$，$X_s$ 或者在代数闭域 $k'=k(s)$ 上平展，或者在 $k'$ 上同构于 $d-2$ 个概形 $\operatorname{Spec}k'$ 与概形
\[
I_{k'}=\operatorname{Spec}\bigl(k'[t]/(t^2)\bigr)
\]
的不交和。\transnote{法稿第60页已写 $\operatorname{Spec}k'$，英译误作 $\operatorname{Spec}k$；正文据法稿恢复。$k'=k(s)$ 沿用底本的几何点记号；这里的 $s$ 是取值于代数闭域的几何点，不能误作底空间任意点的未代数闭剩余域。}

用前辈们的语言来说：每个这样的几何纤维中至多有一个分歧点，而且它是“二次的”。\transnote{法稿第60页明确写“au-dessus de $s$”（在 $s$ 上方）；英译漏掉此范围。正文恢复逐纤维的限制；不能说整个态射总共只有一个分歧点。“二次的”按上面的局部纤维 $\operatorname{Spec}(k'[t]/(t^2))$ 理解；不额外断言该分歧在任意特征下都是温和分歧。这里“射影曲线”明示对原文“$P$ 中曲线”的闭子概形解释：若仅为局部闭曲线，删去曲线上一点即可破坏各纤维恒定长度的结论。原稿在此未另给证明，所用9.2对8.8的依赖仍须保留，不因本轮恢复了不等号和剩余域就声称已补齐全部存在性论证。}

